\oldpage[IV-2]{46}
\subsection{\texorpdfstring{Propiedades de los haces $\mathcal F/t\mathcal F$.}{Propiedades de los haces F/tF}}

\begin{proposition}[3.4.1]
\label{IV.2.3.4.1-fr}
Sean $X$ un preesquema localmente noetheriano, $t$ una sección de
$\mathcal O_X$ sobre $X$, $Y$ el subpreesquema cerrado de $X$ definido por
el ideal $t\mathcal O_X$ de $\mathcal O_X$. Sean $\mathcal F$ un
$\mathcal O_X$-módulo coherente, $S$ el subpreesquema cerrado reducido de
$X$ cuyo espacio subyacente es $\operatorname{Supp}(\mathcal F)$, $(S_i)$ la
familia de los subpreesquemas cerrados reducidos de $X$ cuyos espacios
subyacentes son las componentes irreducibles de $S$; se denota por $s_i$ el
punto genérico de $S_i$. Por último, sean $Z$ una componente irreducible de
$\operatorname{Supp}(\mathcal F/t\mathcal F)=S\cap Y$, $z$ su punto
genérico.
\begin{enumerate}
\item[\textup{(i)}] Para todo $i$ tal que $Z\subset S_i$, $Z$ es una componente
irreducible de $S_i\cap Y$.
\item[\textup{(ii)}] Si $Z$ no es igual a uno de los $S_i$, se tiene
\end{enumerate}
\end{proposition}

\begin{env}[3.4.1.1]
\label{IV.2.3.4.1.1-fr}
\[
\operatorname{long}((\mathcal F/t\mathcal F)_z)
\geqslant\sum_i\operatorname{long}(\mathcal F_{s_i})
\]
\end{env}

\textit{donde la suma del segundo miembro se extiende a todos los $i$ tales que
$Z\subset S_i$.}

\begin{enumerate}
\item[(iii)] \textit{Supongamos que $Z$ no sea igual a ninguno de los $S_i$. Para que
los dos miembros de} (3.4.1.1) \textit{sean iguales, es necesario y suficiente que se
satisfagan las dos condiciones siguientes~:}
\begin{enumerate}
\item[$\alpha$)] \textit{$t_z$ es $\mathcal F_z$-regular} $(\mathbf{0},15.1.4)$.
\item[$\beta$)] \textit{Para todo $i$ tal que $Z\subset S_i$, la imagen canónica del
germen $t_z$ en $\mathcal O_{S_i,z}$ genera el ideal maximal de este anillo}
(lo que implica que $\mathcal O_{S_i,z}$ es un \emph{anillo de valoración discreta} y
que la imagen de $t_z$ es un \emph{uniformizante} de este anillo).
\end{enumerate}
\end{enumerate}

\begin{proof}
(i) Si $j:Y\to X$ es la inyección canónica, se tiene
$\mathcal F/t\mathcal F=\mathcal F\otimes_{\mathcal O_X}\mathcal O_Y=j^*(\mathcal F)$, por tanto
$\operatorname{Supp}(\mathcal F/t\mathcal F)=j^{-1}(S)=S\cap Y$ (\textbf{I}, 9.1.13), de donde se sigue la afirmación.

(ii) y (iii). Como los $s_i$ tales que $Z\subset S_i$ son los que
pertenecen a $\operatorname{Spec}(\mathcal O_z)$, para demostrar (ii) y (iii)
se puede reemplazar $X$ por $\operatorname{Spec}(\mathcal O_z)$; y si
$M=\mathcal F_z$, se puede por tanto suponer que $\mathcal F=\widetilde M$,
de donde $\mathcal F_{s_i}=M_{\mathfrak p_i}$, denotando por $\mathfrak p_i$
los ideales primos mínimos de $\mathcal O_z$. Por otra parte, como $M$ es un
$\mathcal O_z$-módulo de tipo finito, se tiene $S=S'_{\mathrm{red}}$, con
$S'=\operatorname{Spec}(\mathcal O_z/\mathfrak a)$, donde $\mathfrak a$ es
el anulador de $M$ en $\mathcal O_z$ $(\mathbf{0}_{\mathrm I},1.7.4)$, y los dos
miembros de (3.4.1.1) conservan los mismos valores, ya se considere $M$ como
$\mathcal O_z$-módulo o como $(\mathcal O_z/\mathfrak a)$-módulo; se ve así
que finalmente se puede reemplazar $X$ por
$S'=\operatorname{Spec}(A)$, donde $A$ es un anillo local noetheriano, siendo
$M$ un $A$-módulo \emph{fiel}; como $Z$ es cerrado en $S$, la hipótesis de que
$Z\neq S_i$ para todos los $i$ significa que $s_i\notin Z$, por tanto
$\dim(A)>0$; por último, decir que $z$ es punto genérico de $Z$, componente
irreducible de $S_i\cap V(t)$, significa que $A/tA$ es de dimensión $0$ (en
otras palabras, es un anillo local \emph{artiniano}). La cuestión se reduce, pues, a demostrar
el enunciado siguiente:
\end{proof}
